\documentclass[12pt]{amsart} 
\usepackage{mathtools}

\input{./latex-preamble/cm-preamble.tex}
\input{./latex-preamble/use-scr.tex}
\input{./latex-preamble/fr-mod.tex}


%\graphicspath{ {./diagrams/} }

\begin{document}

{Hartshorne} \hfill {9-27-2026}

\bigskip\bigskip

\begin{center}

{\sc Chapter III.9}

\end{center}

\begin{center}

  {Noah Parker} 
    
\end{center}

\bigskip\bigskip


\begin{enumerate}
  \item[II.3.16.] {
      \emph{Noetherian Induction.}
      Let $X$ be a noetherian topological space, and let $\P$ be a property of closed subsets of $X$.
      Assume that for any closed subset $Y$ of $X$, if $\P$ holds for every proper closed subset of $Y$, then $\P$ holds for $Y$. 
      (In particular, $\P$ must hold for the empty set.)
      Then $\P$ holds for $X$.
    }
\end{enumerate}

\begin{proof}
  Noetherianity allows us to write $X=\bigcup_{i=1}^s X_i$ for $X_i$ irreducible, with some integer $n=\sum_{i=1}^s \dim X_i$.
  We induct on $(n,s)$ with the lexicographic ordering.
  If $n=s=0$, then $X=\0$, so $\P$ holds for $X$.
  If $n=0$ and $s>0$, then $X = \{x_1,\ldots, x_s\}$ is the union of finitely many closed points.
  Then every proper subset of $X'\subset S$ is a union of $s'<s$ closed points, so that $\P$ holds for $X'$ by induction, whence $\P$ holds for $X$.

  Now, write $X = \bigcup_{i=1}^s X_i$, reordering so that $\dim X_i>0$ for $1\leqs i\leqs r$ and $X_j = \{x_j\}$ are closed points for $r<j\leqs s$.
  Take a proper closed subset $X'\subsetneq X$.
  We have an irreducible decomposition
  \[
    X' = \bigcup_{i=1}^{s} X'\cap X_i
  \] 
  for $X'$, with some $X'\cap X_i$ empty.
  Let $s'$ be the number of nonempty $X'\cap X_i$, and $n'=\sum_{i=1}^s \dim X'\cap X_i$.

  Take some $i$ such that $X_i\cap X'\subsetneq X_i$ is a proper closed subset.
  If $i>r$, then $X_i\cap X'=\0$ gives us $s'<s$, so $\P$ holds for $X'$ by induction.
  Otherwise we have $\dim X_i\cap X'<\dim X_i$, so that $n'<n$, whence $\P$ holds for $X'$ by induction once more.
  We conclude that, in fact, $\P$ holds for $X$
\end{proof}

\begin{enumerate}
  \item[II.3.17] {
      \emph{Zariski Spaces}.
      A topological space $X$ is a \emph{Zariski space} if it is noetherian and every (nonempty) closed irreducible subset of $X$ has a unique generic point.
      \begin{enumerate}
        \item Show that if $X$ is a noetherian scheme, then $\Sp(X)$ is a Zariski space.
        \item show that any minimal nonempty closed subset of a Zariski space onsists of one point.
  We call these \emph{closed points}.
        \item Show that a Zariski space $X$ satisfies the axiom $T_0$: given any two distinct points of $X$, there is an open set containing one but not the other.
        \item If $X$ is an irreducible Zariski space, then its generic point is contained in every nonempty open subset of $X$.
        \item If $x_0,x_1$ are points of a topological space $X$ ,and if $x_0\in \{x_1\}^-$, then we say $x_1$ \emph{specializes} to $x_0$, written $x_1\rightsquigarrow x_0$.
          We also say $x_0$ is a \emph{specialization} of $x_1$, or that $x_1$ is a \emph{generalization of $x_0$}.
          Now, let $X$ be a Zariski space.
          Show that the minimal points, for the partial ordering determined by $x_1>x_0$ if $x_1\specto x_0$, are the closed points, and the maximal points are the generic points of the irreducible components of $X$.
          Show also that a closed subset contains every specialization of any of its points. (We say closed subsets are \emph{stable under specialization}.)
          Similarly, open subsets are \emph{stable under generalization.}
        \item Show that the maximal spectrum of a ring $A$ is Zariski if and only if it is equal to $\Spec A$.
      \end{enumerate}
    }
\end{enumerate}

\begin{proof}
  (a) Trivial, given (II, Ex. 2.9).

  (b) Let $X$ be a Zariski space, and let $\0\neq Y\subset X$ be a minimal nonempty closed subset of $X$.
  In particular, $Y$ is irreducible, and thus contains a unique generic point $\zeta\in Y$.
  Minimality implies $\{y\}^-= Y$ for all $y\in Y$, but then uniqueness implies $Y=\{\zeta\}$, as desired.

  (c) Let $x,y\in X$ be points such that $x\in U$ open implies $y\in U$.
  Then $y\in C$ closed implies $x\in C$, so that $\{y\in C\textup{ closed}\}\subset \{x\in C\textup{ closed}\}$.
  Hence,
  \[
    \{x\}^-=\bigcap_{x\in C}C\subset \bigcap_{y\in C}C = \{y\}^-.
  \] 
  If $x,y\in X$ have $x\in U$ if and only if $y\in U$, then $\{x\}^-=\{y\}^-$ implies $x=y$.

  (d) If $\0\neq U\subset X$ is a nonempty open set, then $x\notin U$ implies $\{x\}^-\subset U^c\subsetneq X$ is proper, and hence $x$ is not generic.

  (e) Let $x\in X$ be minimal.
  $y\in \{x\}^-$ implies $x\specto y$, so that $y=x$.
  Hence, $x$ is closed.
  Now, let $\zeta\in X$ be maximal.
  $\{\zeta\}^-\subset \{y\}^-$ implies $y\specto \zeta$, so that $y=\zeta$.
  Each irreducible closed subset of $X$ has a generic point by definition, so we have that $\{\zeta\}^-$ is indeed an irreducible component of $X$.

  Let $x\in C$ be contained in a closed set $C$, and take $y$ such that $x\specto y$.
  Then $y\in \{x\}^-\subset C$, as desired.
  Now, we claim that $S\in X$ is closed under generalization if and only if $S^c$ is closed under specialization.
  If $x\in S$ and $y\specto x$, then $y\in S^c$ and $S^c$ closed under specialization would imply $x\in S^c$, a contradiction.
  The converse is identical, and our result for open sets follows.

  (f) One direction is clear.
\end{proof}

\begin{enumerate}
  \item[II.3.18] {
      \emph{Constructible Sets}.
      Let $X$ be a Zariski topological space.
      A \emph{constructible subset} is a subset which belongs to the smallest family $\fF$ such that (1) every open subset is in $\fF$, (2) a finite intersection of elements of $\fF$ is in $\fF$, and (3) the complement of any element of $\fF$ is in $\fF$.
      \begin{enumerate}
        \item A subset of $X$ is \emph{locally closed} if it is the intersection of an open subset with a closed subset.
          Show that a subset of $X$ is constructible if and only if it can be written a finite disjoint union of locally closed subsets.
        \item Show that a constructible subset of an irreducible zariski space $X$ is dense if and only if it contains the generic point.
          Furthermore, in that case it contains a nonempty open subset.
        \item A subset $S$ of $X$ is closed if and only if it is constructible and stable under specialization.
          Similarly, a subset $T$ of $X$ is open if and only if it is constructible and stable under generalization.
        \item If $f:X\to Y$ is a continuous map of Zariski spaces, then the inverse image of any constructible set of $Y$ is constructible in $X$.
      \end{enumerate}
    }
\end{enumerate}

\begin{proof}
  (a) Write $\fL$ for the family of finite disjoint unions of locally closed subsets. 
  It is clear that $\fL\subset \fF$, so it suffices to show that $\fF\subset \fL$.
  We will do so by showing $\fL$ satisfies (1-3), and conclude our result via minimality.

  (1) is trivial.
  If $X=X_1\cap X_2$ for $X_i\in \fL$, then we can write
  \[
    X_i= U_{i1}\cap C_{i1}\sqcup\cdots\sqcup U_{im}\cap C_{im}\quad (i=1,2),
  \]
  adding null sets if necessary so that each is a union of $m$ locally closed sets.
  Then
  \begin{align*}
    X
      &= \left(\bigsqcup_{j=1}^m (U_{1j}\cap C_{1j})\right)\cap \left(\bigsqcup_{j=1}^m (U_{2j}\cap C_{2j})\right) \\
      &= \bigsqcup_{1\leqs j,k\leqs m}(U_{1j}\cap U_{2k})\cap (C_{1j}\cap C_{2k})\in \fL,
  \end{align*}
  Since the union is clearly disjoint.
  Arbitrary finite intersections follow by induction, giving us (2).

  Suppose 
  \[
    X= U_1\cap C_1\sqcup\cdots\sqcup U_m\cap C_m\in \fL.
  \]
  Then
  \begin{align*}
    X^c &= U_1\cap C_1\sqcup\cdots\sqcup U_m\cap C_m \\
        &= (U_1^c\cup C_1^c) \cap\cdots \cap (U_m^c\cup C_m^c) \\
        &= ((U_1\cap C_1^c)\sqcup U_1^c)\cap\cdots \cap((U_m\cap C_m^c)\sqcup U_m^c) \in \fL
  \end{align*}
  by (2), whence we $\fL$ satisfies (3).
  We conclude that in fact $\fF=\fL$

  (b) Let $X$ be an irreducible Zariski space with generic point $\zeta$, and let $Y\in \fF$ be constructible and dense, i.e. $Y^-=X$.
  By (a), we can write $Y=\bigsqcup_{i=1}^n(U_i\cap C_i)\in \fF$ for $U_i\neq \0$ open and $C_i\neq \0$ closed.
  Then certainly $X=Y^-\subset \bigcup_{i=1}^n C_i$, so $X$ irreducible implies $X= C_i$ for some $i$, say $X= C_1$.
  In particular, $\zeta\in C_1$.
  Further, we have $\zeta\in U_1$, as otherwise $\zeta\in U_1^c\subsetneq X$ contradicts genericity of $\zeta$.
  We conclude that, $\zeta\in U_1=C_1\cap U_1\subset Y$, whence both results follow.

  (c) (a) and (Ex. 3.17e) give us one direction, while the proof of the latter gives us that $S$ is closed under specialization if and only if $S^c$ is closed under generalization.
  Now, let $S\in \fF$ be constructible and closed under specialization.
  It suffices to show that $S$ is closed.
  Suppose first that $S=U\cap C\neq \0$ for $C$ irreducible with generic point $\zeta$.
  Then (b) implies $\zeta\in S$, whence $S=C$ by stability under generalization.

  Now, take $S=\bigcup_{i=1}^n U\cap C_i$ for $C_i$ irreducible and $U\cap C_i\neq \0$.
  We claim that each $U\cap C_i$ is closed under specialization.
  Let $x\in U\cap C_i$ and take $x\specto y$.
  Then $y\in C_i$ and $y\in U\cap C_j\subset U$ for some $j$, so, indeed, $y\in U\cap C_i$.
  It follows that the result holds when $S=U\cap C$ for $U$ open and $C$ closed.

  Finally, (a) allows us to write $S=\bigsqcup_{i=1}^n(U_i\cap C_i)$ for $U_i$ open and $C_i$ closed.
  We claim that each $U_i\cap C_i$ is closed under specialization, whence the result would follow from our previous claim.
  Let $x\in U_i\cap C_i$ and take $x\specto y$.
  Then $y\in U_j\cap C_j$ for some $j$, so closure of open sets under generalization gives us $x\in U_j$, 
  Suppose, by way of contradiction, that $i\neq j$.
  Disjointness implies $y\in U_i^c$ and $x\in C_j^c$.
  Then 
  \[
    \{y\}^-\subset \{x\}^- \cap C_j\subset \{x\}^-\setminus \{x\}
  \]
  FINISH THIS!!
  
  (d) If $S=\bigsqcup_{i=1}^n(U_i\cap C_i)$ is constructible in $Y$, then
  \begin{align*}
    f^{-1}(S) &= \bigcup_{i=1}^n f^{-1}(U_i\cap C_i) \\
              &= \left(\bigcap_{i=1}^n (f^{-1}(U_i)\cap f^{-1}(C_i))^c\right)^c
  \end{align*}
  is constructible in $X$.
\end{proof}

\begin{enumerate}
  \item[II.3.19] {
      Let $f:X\to Y$ be a morphism of finite type of noetherian schemes.
      Then the image of any constructible subset of $X$ is a constructible subset of $Y$.
      Prove this theorem in the following steps.
      \begin{enumerate}
        \item Reduce to showing that $f(X)$ itself is constructible, in the case where $X$ and $Y$ are affine, integral noetherian schemes, and $f$ is a dominant morphism.
        \item In that case, show that $f(X)$ contains a nonempty open subset of $Y$ by using the following result from commutative algebra:
          let $A\subset B$ be an inclusion of noetherian integral domains, with $B$ a finitely generated $A$-algebra.
          Then, given a nonzero element $b\in B$, there is a nonzero element $a\in A$ with the following property: if $\phi:A\to K$ is any homomorphism of $A$ to any algebraically closed field $K$, such that $\phi(a)\neq 0$, then $\phi$ extends to a homomorphism $\phi'$ of $B$ into $K$, such that $\phi'(b)\neq 0$.
          [\emph{Hint}: Prove the algebraic result by induction on the number of generators of $B$ over $A$, with a direct proof for $B$ cyclic over $A$. In the application, take $b=1$.]
        \item Now use noetherian induction on $Y$ to complete the proof.
        \item Give some examples of morphisms $f:X\to Y$ of varieties over an algebraically closed field $k$, to show that $f(X)$ need not be either open or closed.
      \end{enumerate}
    }
\end{enumerate}

\begin{proof}
  (a)
\end{proof}

\begin{enumerate}
  \item[9.1.] {
      A flat morphism $f:X\to Y$ of finite type of noetherian schemes is open, i.e. for every open subset $U\subset X$, $f(U)$ is open in $Y$.
      [\emph{Hint}: show that $f(U)$ is constructible and stable under generalization (II, Ex. 3.18) and (II, Ex. 3.19).]
    }
\end{enumerate}

\begin{proof}
  $f(U)$ is constructible by (II, Ex. 3.19), so it suffices to show it is stable under generalization by (II, Ex. 3.18c).
  Suppose $y'\specto y=f(x)$ for some $x\in X$, so $y'$ is the image of some prime $\q'\in\Spec \O_y$ under the inclusion $\Spec \O_y\inc Y$.
  The map $f^\sharp:\O_y\to \O_x$ is flat, so (AM 3.18) gives us that it is in fact surjective.
  \[
    \begin{tikzcd}
      X\ar[r,"f"]&Y \\
      \Spec\O_x\ar[u,hookrightarrow]\ar[r,"f",twoheadrightarrow] &\Spec\O_y\ar[u,hookrightarrow]
    \end{tikzcd}
  \] 
  In particular, it supplies us with a prime $\p'\in \Spec\O_x$ such that $f^\sharp(\p')=\q'$.
  The image $x'\in X$ of $\p'$ gives us $y'=f(x')\in f(X)$, as desired.
\end{proof}

\begin{enumerate}
  \item[9.3.] {
      Some examples of flatness and nonflatness.
      \begin{enumerate}
        \item If $f:X\to Y$ is a finite surjective morphism of nonsingular varieties over an algebraically closed field $k$, then $k$ is flat.
        \item Let $X$ be a union of two planes meeting at a point, each of which maps isomorphically to a plane $Y$.
          Show that $f$ is not flat.
          For example, let $Y=\Spec k[x,y]$ and $X=\Spec[x,y,z,w]/(z,w)\cap (x+z,y+w)$.
        \item Again let $Y=\Spec k[x,y]$, but take $X=\Spec k[x,y,z,w]/(z^2, zw, w^2, xz-yw)$.
          Show that $X_\red\cong Y$ has no embedded points, but that $f$ is not flat.
      \end{enumerate}
    }
\end{enumerate}

\begin{enumerate}
  \item[9.4.] {
      \emph{Open Nature of Flatness}.
      Let $f:X\to Y$ be a morphism of finite type of noetherian schemes.
      Then $\{x\in X|\textup{$f$ if flat at $x$}\}$ is an open subset of $X$ (possibly empty).
    }
\end{enumerate}

\begin{proof}
  
\end{proof}

\begin{enumerate}
  \item[9.5.] {
    }
\end{enumerate}

\begin{enumerate}
  \item[9.7.] {
      Let $Y\subset X$ be a closed subscheme, where $X$ is a scheme of finite type over a field $k$.
      Let $D=k[t]/t^2$ be the ring of dual numbers, and define an \emph{infinitesimal deformation of $Y$ as a closed subscheme of $X$}, to be a closed subscheme $Y'\subset X\times_k D$, which is flat over $D$ and whose closed fibre is $Y$.
      Show that they are classified by $H^0(Y,\N_{Y/X})$, where
      \[
        \N_{Y/X} = \sHom_{\O_X}(\I_Y/\I_Y^2,\O_Y).
      \] 
    }
\end{enumerate}

\begin{proof}
  We consider first the affine case, so that $X=\Spec A$ and $Y=\Spec A/I$.
  Then $X\times_k D=\Spec B$ with $A[t]/t^2$, and $Y'=\Spec B/I'$ for 
  \[
    I'=(a_1 + b_1t,\ldots, a_r +b_rt,b_{r+1}t,\ldots, b_st)
  \]
  with $b_1,\ldots, b_s\in A$ and $(a_1,\ldots, a_r)=I$.
  Hence, the coordinate ring of the closed fibre is given by
  \begin{align*}
    \frac{B}{I'}\otimes_{k[t]/t^2}\frac{k[t]/t^2}{(t)}
     &\cong \frac{B/(t)}{I'/(t)}
     \cong \frac{A}{I}.
  \end{align*}
  Note that any $I'$ can be written in this way without duplicate $a_i$, as $a_i + b_i't = a_i +b_it + (b_i'-b_i)t$.
  It suffices to find a condition for $Y'\to D$ to be flat.

  Recall that a morphism of affine schemes is flat if and only if the associate map of rings is flat.
  (9.1Aa) gives us that $k[t]/t^2\to B/I'$ is flat if and only if the multiplication map $\a\otimes B/I'\to B/I'$ is injective for every finitely generated ideal $\a\lideal k[t]/t^2$.
  There are only two ideals of $k[t]/t^2$, and $(1)\otimes B/I'\to B/I'$ is an isomorphism.
  Thus, we need only consider $(t)\otimes B/I'\to B/I'$.

  First, we note that $t\otimes bt = 0$.
  Then if $b\in B/I'\setminus (t)/I'$

  Now, consider  $\N_{Y/X}=\Hom_A(I/I^2,A/I)$

  The universal property of $X\times_k D$ gives us that maps $Y'\to X\times_k D$ are in correspondence with commuting diagrams of the following form.
  \[
    \begin{tikzcd}[column sep = 0.3cm]
      &X\ar[r] & \Spec k \\
      &X\times_k D \ar[u]\ar[r] & D\ar[u] \\
      Y'\ar[ur,dotted,"\exists!"']\ar[uur, bend left = 30]\ar[urr, bend right = 30]
    \end{tikzcd}
  \] 
\end{proof}

\begin{enumerate}
  \item[9.10.] {
      A scheme $X_0$ over a field $k$ is \emph{rigid} if it has no infinitesimal deformations.
      \begin{enumerate}
        \item Show that $\PP_k^1$ is rigid, using (9.13.2).
      \end{enumerate}
    }
\end{enumerate}

\begin{proof}
  It suffices via (9.13.2) to show that $H^1(X,\T_X)=0$, with $X=\PP_k^1$.
  \begin{align*}
    \T_X &=\Omega_X^\vee \\
         &=\omega_X^\vee \\
         &\cong \O(2).
  \end{align*}
  Then (7.1) and (II, 5.1) give us
  \begin{align*}
    H^1(X,\T_X)^\vee &\cong \Hom(\T_X,\omega_X) \\
                     &\cong \Gamma(X,\O(-4)) \\
                     &= 0,
  \end{align*}
  since $-4<0$.
\end{proof}

\end{document}
